\section{Related Work}
\label{sec:related}

\subsection{County preservation as a legal criterion}

The preservation of political subdivisions has a long history in
redistricting law.  \emph{Reynolds v.\ Sims}~\citep{reynolds1964sims}
acknowledged that states may consider county boundaries as a secondary
criterion after population equality.  The National Conference of State
Legislatures documents subdivision-preservation requirements in 34 states,
with Iowa's constitutional provision being among the most stringent.

Grofman et al.~\citep{grofman1990representation} reviewed the history of
county preservation as a redistricting criterion, noting that courts have
generally upheld county preservation as a valid, non-partisan criterion
that can justify minor deviations from strict population equality.

Benton~\citep{benton2002counties} documented the service-delivery role
of counties --- as units of government responsible for health, courts,
and elections --- as the substantive reason for preserving county
integrity in legislative maps.

\subsection{Algorithmic approaches to subdivision preservation}

Altman~\citep{altman1998compactness} proposed a bicriterion optimisation
approach that treats compactness and subdivision preservation as competing
objectives, finding that small compactness losses can achieve large
subdivision-preservation gains.

Ricca et al.~\citep{ricca2008local} implemented tabu search with county
preservation as a hard constraint, finding that infeasibility arises for
approximately 15--20\% of their test instances --- consistent with our
hard-constraint analysis.

The GerryChain framework~\citep{deford2019recombination} supports optional
county-preservation filters in the ReCom chain, but these function as
post-hoc filters (rejecting maps that split too many counties) rather than
as an edge-weighting signal, making them less efficient for large instances.

\subsection{Edge weighting for graph partitioning}

Edge weighting in graph partitioning is well-studied for non-geographic
applications.  Karypis and Kumar~\citep{karypis1998} showed that METIS
supports general edge weights without loss of correctness or the $O(m)$
per-level complexity.

Bozkaya et al.~\citep{bozkaya2003zone} used a form of edge weighting in
their tabu search for political districting, penalising inter-county
edges with a multiplier.  Their multiplier was fixed at 2.0 based on expert
judgment; our paper provides the first systematic empirical calibration
of this parameter.

\subsection{Multi-criterion redistricting}

Swamy et al.~\citep{swamy2011least} formulated redistricting as a
least-squares multi-criterion problem.  County preservation was one of
several criteria, but they used a hard-constraint formulation that scaled
poorly beyond $n \approx 500$ tracts.

The soft-constraint approach of this paper is related to penalty methods in
constrained optimisation~\citep{mehrotra1998optimal}: the county-sticky
multiplier $\alpha_c$ functions as a penalty on inter-county cuts, with
the penalty calibrated to achieve a target constraint satisfaction rate.
